\documentclass[
  a4paper,
  12pt,
  oneside,
  parskip=half,
  listof=totoc,
  bibliography=totoc,
  DIV=12,
  numbers=noenddot,
  BA.tex
  ]{subfiles}

\begin{document}

\section{Summary and Outlook}

\subsection{ALICE Run 3 measurements}

The jet-induced diffusion wake was studied using the correlation of soft hadrons with a dijet/dihadron trigger pair in ALICE Run 3 0-10\% Pb--Pb data at $\sqrt{s_{NN}} = 5.36\bunit{TeV}$. No correction for detector acceptance or tracking efficiency effects was made.

The analysis results using a correlated dijet pair as trigger revealed anomalies, which are likely to be systematic errors, at $|\eta| = 0.3$, especially for lower jet trigger $p_T$ thresholds. Apart from that, a signal that has the characteristics of a diffusion wake could be identified among all trigger variants, \newline $(40,20)\bunit{GeV/c}, (30,15)\bunit{GeV/c} \text{, and } (20,10)\bunit{GeV/c}$.

The result of the $(40,20)\bunit{GeV/c}$ trigger variant was compared to another ALICE Run 3 analysis that uses 0-30\% Pb--Pb data. This 0-30\% data includes one Pb--Pb dataset that is also used in the 0-10\% analysis. The 0-10\% result is similar to the 0-30\% result, although the two results are slightly shifted relative to each other. The 0-30\% signal is symmetric in pseudorapidity, the 0-10\% data is slightly asymmetric. The 0-30\% data shows a smaller residual signal for higher soft hadron bins, as expected from the predicted diffusion wake signal. The remaining residual is likely to come from systematic errors. The 0-10\% result for the $(40,20)\bunit{GeV/c}$ trigger variant was also directly compared to a CoLBT-hydro prediction \cite{zhong_yang_colbt-hydro_2025}. The prediction provides data for the soft hadron bins $1 \leq p_T \leq 2 \bunit{GeV/c}$ and $2 \leq p_T \leq 4 \bunit{GeV/c}$. The $1 \leq p_T \leq 2 \bunit{GeV/c}$ prediction fits the 0-10\% result well for $\eta > 0$, but not for $\eta < 0$. There, the data is slightly off the prediction's bounds. The $2 \leq p_T \leq 4 \bunit{GeV/c}$ prediction fits the ALICE 0-10\% data well.

To avoid bias effects that come primarily from correlations between the event plane and reconstructed jets, the analysis was also done with a dihadron pair as trigger. The results of the dihadron analysis are not subject to the anomalies found in the dijet analysis at $|\eta| = 0.3$. Furthermore, the statistical errors are smaller than in the dijet analysis across all trigger variants, $(30,15)\bunit{GeV/c}, (20,15)\bunit{GeV/c}, (15,10)\bunit{GeV/c} \text{, and } (15,8)\bunit{GeV/c}$. The $(15,8)\bunit{GeV/c}$ trigger variant was added afterwards because the $(15,10)\bunit{GeV/c}$ variant showed promising results. A signal with diffusion wake characteristics is clearly visible in many $p_T$ bins. The $(15,8)\bunit{GeV/c}$ trigger variant for soft hadrons with transverse momenta $0 \leq p_T \leq 2 \bunit{GeV/c}$ holds the clearest results in terms of uncertainties and symmetry of the diffusion wake-like signal.

To refine the diffusion wake measurement in ALICE data, acceptance and efficiency corrections need to be done. Those studies were out of scope for this thesis.

\subsection{Simulated Pb--Pb event measurements}

The jet-induced diffusion wake was also studied using the correlation of soft hadrons with a dijet/dihadron trigger pair in artificial Pb--Pb data derived from real Pb--Pb data and the PYTHIA event generator.

As part of the studies of bias effects, a simulation was performed to estimate how the jet position is modified by flow modulated bulk particles. The results are presented separately for different jet transverse momentum bins. The result shows that a reconstructed jet's azimuthal position is biased towards the event plane. The effect is strongest for the softest jets. For reconstructed jets with transverse momenta of $20 \leq p_T \leq 40 \bunit{GeV/c}$, the shift is highest when the jet is approximately $d = \pm \pi/4 \bunit{rad}$ away from the event plane. There, the jet is shifted by $0.836 \bunit{deg}$ towards the event plane. For higher jet $p_T$ the value decreases rapidly. The introduced low-particle model of the cluster algorithm predicts the largest shift when $d = \pm \pi/4 \bunit{rad}$ exactly, in the measurement the maximum is slightly off. The model uses a few approximations whose validity may need further attention in order to solve the slight deviation of the model from the measurement.

The study of the diffusion wake analysis with simulated Pb--Pb events, where no diffusion wake is modeled, showed a significant difference signal that opposes the predicted diffusion wake signal by CoLBT-hydro \cite{zhong_yang_colbt-hydro_2025}. The opposing signal is smaller ($\approx 1/8$) than the diffusion wake signal, but of same shape albeit inverted in $\eta$. The signal is significant with $\chi^2 > 20$ for the $(20,10)\bunit{GeV/c}$ dijet trigger and smaller, but non-zero ($\chi^2 = 4.03$) for the $(15,10)\bunit{GeV/c}$ dihadron trigger. The absolute value of the signal is stronger when using a dijet trigger compared to when using a dihadron trigger. The bias increases with decreasing dijet/dihadron trigger pair $p_T$ threshold. 

The bias signal could originate in anisotropic hadron flow and, directly related, in the event selection. The bias signal was studied by magnifying the simulated directed flow low-$p_T$ hadrons are subject to. The, by a factor of 100, increased directed flow revealed a difference signal in the dihadron analysis that has similar characteristics as the CoLBT-hydro diffusion wake signal \cite{zhong_yang_colbt-hydro_2025}. A mechanism that explains the dependence of the signal on the strength of the directed flow was introduced in the case of a dijet trigger. This mechanism may also be valid for a dihadron trigger, but this needs further study on the anisotropic flow of high-$p_T$ hadrons. Studies on this are ongoing within the ALICE collaboration. It needs further studies and more parameter manipulation to determine which parameter change could generate the predicted diffusion wake signal magnitudes. The elliptic and triangular flow are also parameters whose influence on the analysis result can be studied in further simulations. It is of great importance to refine the simulation. Specifically, to implement flow effects for a variety of kinematics for one specific collision energy $\sqrt{s_{NN}}$ and centrality category. To be precise, the flow parameters $v_n$ depending on pseudorapidity $\eta$ and transverse momentum $p_T$ must be known in a big range, e. g. $0 \leq p_T \leq 40 \bunit{GeV}$ to make quantitatively reliable statements on such biases. In this early stage of the simulation, a broad variety of flow coefficient measurements has been used. This includes several centrality categories, different center of mass energies and small $p_T$ windows in which the measurements are valid in. Nevertheless this analysis has shown that not only the diffusion wake could generate a signal of such kind. Further studies need to be conducted to ensure that the diffusion wake-like signal is not part of a systematic bias due to flow. 


\end{document}